\documentclass[10pt]{article}
\usepackage[T1]{fontenc}
\usepackage{lmodern,amsmath,amssymb,amsthm,microtype}
\usepackage[margin=.82in]{geometry}
\usepackage[colorlinks=true,urlcolor=blue]{hyperref}
\newtheorem{theorem}{Theorem}
\newtheorem{lemma}{Lemma}
\newtheorem{corollary}{Corollary}
\title{Explicit Modular Periods for Ordered Partitions into Odd Blocks}
\author{Qichao Wang\\{\small Hebei University of Technology}\\[.3em]
Daoyu Dong\\{\small University of Electronic Science and Technology of China}}
\date{October 27, 2026}
\begin{document}\maketitle
\begin{abstract}
The number of ordered set partitions into odd blocks is eventually
periodic modulo every positive integer. We prove the conjecture in
OEIS A006154 with an explicit period dividing twice the Euler
totient of the modulus; for odd moduli the factor two is unnecessary.
A differential-operator congruence treats all prime powers uniformly.
The least eventual period modulo ten is four, beginning at index one.
\end{abstract}

Let $a(n)$ count ordered partitions of an $n$-element set into blocks
of odd cardinality, with $a(0)=1$. Choosing the first block gives
\[
 a(n)=\sum_{\substack{1\le j\le n\\ j\text{ odd}}}
       \binom nj a(n-j),\qquad
 F(x)=\sum_{n\ge0}a(n)\frac{x^n}{n!}=\frac1{1-\sinh x}.
\]
OEIS A006154 records this sequence and Peter Bala's April 12, 2023
conjecture that its reduction modulo every positive integer is
eventually periodic~\cite{oeis}.

\begin{theorem}\label{thm:main}
Factor $m\ge2$ into prime powers. Define
\[
 P(m)=\operatorname{lcm}\left(
 \{\varphi(p^e):p^e\parallel m,\ p\text{ odd}\}
 \cup\{2^e:2^e\parallel m\}\right).
\]
The second set is empty for odd $m$. For every $n\ge m$,
\[
 a(n+P(m))\equiv a(n)\pmod m.
\]
In particular $P(m)\mid2\varphi(m)$; for odd $m$,
$P(m)\mid\varphi(m)$. Modulo one the sequence is constant.
\end{theorem}

\section{Differential congruences}
Work in the torsion-free ring
$A=\mathbb Z[s,c,(1-s)^{-1}]$, and put $u=(1-s)^{-1}$.
The derivation
\[
 D=c\partial_s+s\partial_c
\]
preserves $A$, since $D(u)=cu^2$. Substitution
$(s,c)=(\sinh x,\cosh x)$ turns $D$ into differentiation with
respect to $x$. Thus, with $\varepsilon(f)=f(0,1)$,
\begin{equation}\label{eq:evaluation}
 a(n)=\varepsilon(D^n u).
\end{equation}
This formula uses only integer expressions and no division by
factorials modulo a modulus.

\begin{lemma}\label{lem:prime}
As endomorphisms of $A$,
\[
 D^p\equiv D\pmod p\quad(p\text{ odd prime}),
 \qquad D^4\equiv D^2\pmod2.
\]
\end{lemma}
\begin{proof}
In characteristic $p$, the $p$th iterate of a derivation is again
a derivation, by the iterated product rule and the divisibility
of $\binom pj$ for $0<j<p$. For odd $p$, the derivations $D^p$
and $D$ agree on $s,c$, since $Ds=c$ and $Dc=s$. They therefore
agree on the polynomial ring and its localization. In
characteristic two both $D^4$ and $D^2$ are derivations; on each
of $s,c$ they equal the identity. They agree on the localization
as well. This proves both operator congruences.
\end{proof}

\begin{lemma}\label{lem:lift}
Let $X,Y$ be commuting integer-linear endomorphisms of a torsion-free
abelian group. If $X\equiv Y\pmod{p^r}$, $r\ge1$, then
$X^p\equiv Y^p\pmod{p^{r+1}}$.
\end{lemma}
\begin{proof}
Write $X=Y+p^rZ$. The endomorphism $Z$ is integer-linear and
commutes with $Y$, by torsion-freeness. In the binomial expansion
of $(Y+p^rZ)^p-Y^p$, the first term has factor $p^{r+1}$.
For $2\le j<p$ the $j$th term has factor $p^{rj+1}$,
and the last has factor $p^{rp}$, where $rp\ge r+1$.
Every term is divisible by $p^{r+1}$.
\end{proof}

\begin{proof}[Proof of Theorem~\ref{thm:main}]
Repeatedly apply Lemma~\ref{lem:lift} to the commuting powers of
$D$ in Lemma~\ref{lem:prime}. This gives
\[
 D^{p^e}\equiv D^{p^{e-1}}\pmod{p^e}\quad(p\text{ odd}),
 \qquad D^{2^{e+1}}\equiv D^{2^e}\pmod{2^e}.
\]
Multiplication by further powers of $D$ and
\eqref{eq:evaluation} prove periods $p^{e-1}(p-1)$ from
index $p^{e-1}$ for odd $p$, and period $2^e$ from index
$2^e$ for the power of two. Every threshold is at most $m$.
The Chinese remainder theorem gives their least common multiple
as a period modulo $m$. Finally each odd-prime period is the
corresponding factor of $\varphi(m)$, and the power-of-two period
is twice that factor, proving the divisibility assertions.
\end{proof}

\section{The exact period modulo ten}
Modulo five, Lemma~\ref{lem:prime} gives period four from index
one. To handle the initial indices modulo two, reduce in
\[
 V_2=(\mathbb Z/2\mathbb Z)[s,c]/(s^2,c^2-1).
\]
The derivation preserves the ideal and $u=1+s$ in this quotient.
Consequently
\[
 Du=c,\qquad D^2u=s,\qquad D^3u=c=Du.
\]
Evaluation at $(0,1)$ descends to $V_2$, proving period two from
index one. Combining these results gives
\[
 a(n)\bmod10=(1,2,7,2)_{\,1+((n-1)\bmod4)}\qquad(n\ge1).
\]
The block has neither period one nor period two, so its least
eventual period is four. This also proves the specific example
recorded with the general conjecture.

The supplementary verifier compares all 201 published terms with
the first-block recurrence and independently enumerates literal
set partitions through eight elements. Finite differential quotients
check the operator calculation for moduli two through 32. The
proof above establishes the universal periodicity; the finite
checks audit its implementation.
\begin{thebibliography}{1}
\bibitem{oeis} The OEIS Foundation, A006154, including P. Bala's
periodicity conjecture (2023) and differential formula (2011),
\url{https://oeis.org/A006154}, accessed October 27, 2026.
\end{thebibliography}\end{document}
